% ===========================================================================
\part{Orientation}
% ===========================================================================

\chapter{How to read this document}

\section{What this document is for}

You have written the skeleton of a personal website in Rust. The database exists, four
migrations have been applied, four domain modules each have a model and a repository, and
the server compiles. What does \emph{not} yet exist is any connection between the data and the
page: \file{src/app.rs} is still the Leptos starter template with a click counter in it.

This document exists so that the next stretch of work --- the part where the data finally
reaches the screen --- is built on decisions made deliberately rather than by accident. That
distinction matters more in this stack than in most. A Leptos application compiles the same
source file twice, once for a native server and once for WebAssembly in the browser, and the
consequences of that fact reach into the type of every struct you write. If you place a struct
in the wrong module, or gate a dependency behind the wrong feature flag, you do not get a
runtime bug you can debug later --- you get a compile error whose message points somewhere
unhelpful, and the natural way out of it is to work around the symptom and entrench the mistake.
Several such mistakes are already latent in the code. They are catalogued here, with the
reasoning that makes them obvious in hindsight.

\begin{keyidea}
The goal is not to describe the system you have. It is to give you the mental model that makes
the right next move feel obvious, and the concrete list of things to fix before that move
becomes harder.
\end{keyidea}

\section{How it is organised}

The document runs from concepts to consequences to a plan.

\begin{description}
  \item[Part I --- Orientation] takes inventory. What is in the repository today, what
    each dependency is actually responsible for, and what the toolchain does when you type
    \code{cargo leptos watch}.

  \item[Part II --- The mental model] is the part to read even if you skip everything else.
    It explains the two-compilation model, the lifecycle of a request through server rendering
    and hydration, and the boundary that separates code allowed into the browser from code that
    must never go there. Almost every structural decision later in the document is a consequence
    of these three things.

  \item[Part III --- Architecture] turns the model into layout: where domain logic lives, why a
    persistence model and a transport DTO must be two different types, why the existing
    \code{Repository<T>} trait will not survive contact with async code, and how the connection
    pool should reach a server function.

  \item[Part IV --- Data] reviews the schema as it is against the schema it should be, and
    establishes migration discipline. Three of the four existing migrations cannot be rolled
    back; this part shows why and how to repair them.

  \item[Part V --- The front end] covers page structure, the data-loading patterns Leptos
    provides, and the exact configuration for Tailwind CSS, which you have said you intend to adopt.

  \item[Part VI --- Quality and operations] describes the testing strategy this stack rewards,
    plus logging, configuration, containerisation and deployment.

  \item[Part VII --- The plan] is the actionable residue: a numbered register of every defect
    found, a milestone-ordered roadmap, and the conventions worth committing to now, while the
    codebase is still small enough that consistency is cheap.
\end{description}

\section{Conventions used here}

Five kinds of marginal box appear throughout. They exist so you can skim for the parts you need.

\begin{pattern}
A pattern to adopt. If you read only these boxes you would still end up with a
defensible architecture.
\end{pattern}

\begin{antipattern}
Something to avoid, usually because it is already present in the code and will
cost more to remove later than now.
\end{antipattern}

\begin{pitfall}
A trap that is not wrong exactly, but that will waste an afternoon if you meet it
without warning.
\end{pitfall}

\begin{why}
The reasoning behind a recommendation. Advice you cannot justify is advice you cannot
adapt when circumstances change, so the justification is given whenever it is not obvious.
\end{why}

\begin{verified}
A claim confirmed by running something on this machine, with the real output quoted.
Where a box like this is absent and a factual claim is made about tool behaviour, treat it as
documented-but-unverified.
\end{verified}

\vspace{2mm}
Defects are numbered \dnum{1}, \dnum{2}, and so on. The numbers are assigned in
Chapter~\ref{ch:register} and referenced from wherever the underlying topic is discussed, so you
can read the register first and jump to the explanation, or read forward and collect the
numbers as you go.

% ===========================================================================
\chapter{The state of the code today}
\label{ch:state}

\section{An honest inventory}

The repository holds 589 lines of Rust across 25 files, plus four migrations, a
\file{docker-compose.yml} and the untouched remains of the Leptos Actix starter template. The
useful way to read that inventory is not file-by-file but by asking, of each file, \emph{is this
load-bearing?} Table~\ref{tab:inventory} does that.

\begin{table}[htbp]
\centering
\caption{Every Rust source file, and whether it currently does anything.}
\label{tab:inventory}
\small
\begin{tabular}{@{}L{0.325\textwidth}rL{0.13\textwidth}L{0.36\textwidth}@{}}
\toprule
\textbf{File} & \textbf{Lines} & \textbf{Status} & \textbf{Note} \\
\midrule
\file{src/main.rs}          & 88 & live & Starter template, unmodified. No pool, no server-fn wiring. \\
\file{src/lib.rs}           & 27 & live & Module roots; every server module gated on \code{ssr}. \\
\file{src/app.rs}           & 66 & live & Starter template. A counter and a 404 page. \\
\file{src/schema.rs}        & 103 & live & Diesel-generated, matches the live database. \\
\midrule
\file{src/core/mod.rs}      & 4 & live & The \code{Repository<T>} trait. Read-only. \\
\file{src/database/connection.rs} & 17 & \textbf{orphan} & Builds a pool that nothing ever calls. \\
\midrule
\file{src/portfolio/model.rs} & 26 & live & \code{PortfolioItem}, \code{Tag}. Public fields. \\
\file{src/portfolio/repository.rs} & 44 & \textbf{orphan} & Compiles; never constructed. \\
\file{src/publications/model.rs} & 14 & live & \code{PublicationItem}. Public fields. \\
\file{src/publications/repository.rs} & 43 & \textbf{orphan} & Compiles; never constructed. \\
\file{src/job\_experience/model.rs} & 26 & live & \textbf{Private fields.} \\
\file{src/job\_experience/repository.rs} & 46 & \textbf{orphan} & Compiles; never constructed. \\
\file{src/personal\_information/model.rs} & 27 & live & \textbf{Private fields.} \\
\file{src/personal\_information/repository.rs} & 47 & \textbf{orphan} & Compiles; never constructed. \\
\midrule
\file{src/entrypoint/mod.rs} & 2 & live & Declares the two files below. \\
\file{src/entrypoint/http.rs} & 0 & \textbf{empty} & Declared and compiled; contains nothing. \\
\file{src/entrypoint/application\_state.rs} & 0 & \textbf{empty} & Same. Created during this
document's writing. \\
\bottomrule
\end{tabular}
\end{table}

\begin{pitfall}[This table has a shelf life]
The \file{src/entrypoint/} rows changed twice while this document was being written. As of writing,
\file{routes/} has been deleted and replaced by \file{http.rs} and
\file{application\_state.rs}, both empty but both properly declared --- which is a genuine
improvement over the previous state, described below. Where this document discusses that directory,
it discusses the shape rather than the file list. Chapter~\ref{ch:composition} is where
\code{application\_state.rs} gets an opinion.
\end{pitfall}

Three words in that table deserve definitions, because the difference between them is the
difference between a gap and a bug.

\begin{description}
  \item[live] The compiler sees it and something depends on it.
  \item[orphan] The compiler sees it, it type-checks, but no code path reaches it. Orphans are
    normal in a half-built system --- they are work in progress. The risk is that an orphan can
    be quietly wrong for weeks, because nothing exercises it.
  \item[dead] The compiler never sees it at all. This is worse than an orphan, because you get
    no type-checking, no warnings, and no signal that the file is not part of the program.
\end{description}

\subsection{Dead modules, and how to detect one}

Until an hour before this was written, \file{src/entrypoint/} contained a subtree the compiler had
never seen. It is worth reconstructing, because the mechanism is invisible unless you trace module
declarations by hand and it will happen again.

Rust does not discover files. A file is part of the crate only if some ancestor module declares it
with \code{mod}. The chain broke immediately:

\begin{lstlisting}[style=rs,caption={The declaration chain as it was, traced.}]
// src/lib.rs
#[cfg(feature = "ssr")]
pub mod entrypoint;            // <-- declares entrypoint, so mod.rs IS read

// src/entrypoint/mod.rs
                               // <-- EMPTY. Declares nothing.
                               //     Therefore http.rs and routes/ were never read.

// src/entrypoint/routes/job_experience/mod.rs   (unreachable anyway)
pub mod request;
pub mod response;              // <-- and note: route.rs was not declared here either
\end{lstlisting}

Five files, four of them empty, none of them compiled --- and, tellingly, even the one file that
did declare its siblings had forgotten \file{route.rs}. Two independent layers of the same mistake.

\begin{verified}[Verified: the compiler really does not see such files]
An investigation for this document wrote the text \code{this is not valid rust !!!} into three of
those files and ran \code{cargo check --features ssr}. Zero errors mentioned \code{entrypoint};
rustc never opened them. The files were then restored.
\end{verified}

\begin{pitfall}[The silence of dead modules]
An undeclared \code{.rs} file produces no warning of any kind. \code{cargo check} is clean,
\code{clippy} is clean, and your editor will happily give you completions inside a file the
compiler is ignoring. If you are ever unsure whether a file is in the build, put a deliberate
syntax error in it and run \code{cargo check}. Silence means it is dead.
\end{pitfall}

The current state is better --- \file{http.rs} and \file{application\_state.rs} are both declared,
so both are compiled --- and it has traded one smell for a milder one: two declared modules that are
empty. That is fine as scaffolding and worth watching, because an empty declared module is a
promise. Chapter~\ref{ch:composition} takes a position on what \code{application\_state.rs} should
contain, since a state struct at the composition root is a good idea rather than a leftover.

\section{What the database looks like}

The database is further along than the Rust. A PostgreSQL 17 container is running, and all four
migrations plus Diesel's bookkeeping migration have been applied to a database named
\code{my-site}.

\begin{verified}
\begin{lstlisting}[style=out]
$ diesel migration list
Migrations:
  [X] 00000000000000_diesel_initial_setup
  [X] 2026-08-02-211314-0000_create_portfolio_item
  [X] 2026-08-03-231439-0000_create_publications_item
  [X] 2026-08-03-232042-0000_create_job_experience
  [X] 2026-08-03-233332-0000_create_personal_information

$ psql -c 'select count(*) from portfolio_items'   ->  0
  ... and 0 rows in every other table.
\end{lstlisting}
\end{verified}

Seven tables exist and every one of them is empty. That is worth knowing before you write the
first page, because a page that renders an empty list looks identical to a page whose query is
broken. Chapter~\ref{ch:seed} covers seeding, and it should happen \emph{before} the first
component, not after.

\begin{pitfall}[Two databases, one of them a decoy]
\file{docker-compose.yml} creates a database called \code{mydatabase}. Your \file{.env} points
\code{DATABASE\_URL} at a database called \code{my-site}. Both exist on the server --- \code{my-site}
because \code{diesel setup} created it, \code{mydatabase} because Compose did. \code{mydatabase}
is empty and unused. Anyone who clones this repository, runs \code{docker compose up}, and trusts
the Compose file will connect to the wrong, empty database and see nothing. This is
defect \dnum{14}.
\end{pitfall}

\section{What is missing, stated plainly}

It is easy to lose the shape of the remaining work in a list of small defects, so here is the
shape. Four things are absent, and they are absent in a specific order of dependency:

\begin{enumerate}
  \item \textbf{No transport.} There is no server function, no route, and no HTTP endpoint that
    returns domain data. The repositories can read from Postgres, and nothing can call the
    repositories.
  \item \textbf{No shared types.} Even if a transport existed, the current models cannot cross
    it. They are compiled only in the server build, and they derive \code{Serialize} from a
    dependency that only the server build enables. Chapter~\ref{ch:boundary} works through
    exactly what happens.
  \item \textbf{No presentation.} \file{src/app.rs} has one route and a counter. There are no
    pages for portfolio, publications, experience or contact, no layout, no navigation.
  \item \textbf{No writes.} \code{Repository<T>} offers \code{get\_all} and \code{get\_by\_id}.
    There is no insert, update or delete anywhere in the codebase, and no path by which you could
    author content other than by hand-writing SQL.
\end{enumerate}

Everything else --- Tailwind, tests, CI, deployment, tracing --- is genuinely optional until those
four are done. Part~VII sequences them.

% ===========================================================================
\chapter{The stack, and who is responsible for what}
\label{ch:stack}

Six moving parts sit between a browser request and a row in Postgres. Confusion about which one
owns a given responsibility is the most common source of misplaced code in this stack, so this
chapter draws the boundaries once.

\begin{diagram}{The stack, with each component's single responsibility. Solid arrows are
runtime data flow; the dashed arrow is a build-time relationship.}{fig:stack}
\begin{tikzpicture}[node distance=6mm]
  \node[cli,minimum width=34mm] (br) {\textbf{Browser}\\[1pt]\scriptsize DOM, fetch, history};
  \node[cli,minimum width=34mm,below=8mm of br] (wasm) {\textbf{wasm bundle}\\[1pt]\scriptsize
    \code{leptos} (hydrate) + your \code{app.rs}};
  \node[srv,minimum width=40mm,below=13mm of wasm] (actix) {\textbf{Actix-web 4}\\[1pt]\scriptsize
    HTTP server, routing, static files};
  \node[srv,minimum width=40mm,below=8mm of actix] (leptos) {\textbf{leptos + leptos\_actix}\\[1pt]\scriptsize
    renders HTML, hosts server functions};
  \node[srv,minimum width=40mm,below=8mm of leptos] (domain) {\textbf{your domain + repositories}\\[1pt]\scriptsize
    business rules, queries};
  \node[srv,minimum width=40mm,below=8mm of domain] (diesel) {\textbf{Diesel 2.3 + r2d2}\\[1pt]\scriptsize
    typed SQL, connection pool};
  \node[db,below=8mm of diesel,minimum width=32mm] (pg) {\textbf{PostgreSQL 17}};

  \node[shared,right=16mm of wasm,minimum width=30mm,align=center] (cl) {\textbf{cargo-leptos}\\[1pt]\scriptsize
    builds both sides,\\ runs sass/tailwind,\\ watches and reloads};

  \draw[flow] (br) -- (wasm);
  \draw[flow] (wasm) -- node[lbl,right]{HTTP} (actix);
  \draw[flow] (actix) -- (leptos);
  \draw[flow] (leptos) -- (domain);
  \draw[flow] (domain) -- (diesel);
  \draw[flow] (diesel) -- (pg);
  \draw[dep] (cl) -- (wasm);
  \draw[dep] (cl.east) -- ++(6mm,0) |- (actix.east);

  \begin{scope}[on background layer]
    \node[lane,fit=(actix)(diesel)(pg),inner sep=4mm,label={[tag,anchor=north east]north east:server process}] {};
    \node[lane,fit=(br)(wasm),inner sep=3.5mm,label={[tag,anchor=north east]north east:client}] {};
  \end{scope}
\end{tikzpicture}
\end{diagram}

\section{The pieces, one at a time}

\subsection*{Leptos 0.8 --- the reactive UI framework}

Leptos is the only piece that exists on both sides of the network. You write components once, as
Rust functions returning \code{impl IntoView}, and Leptos either renders them to an HTML string
(on the server) or mounts them onto existing DOM and makes them interactive (in the browser).
Which of those two things it does is chosen by a Cargo feature, not at runtime. That single fact
is the source of most of Part~II.

Leptos also owns two things people often expect Actix to own. \emph{Routing of application pages}
is a Leptos concern --- \code{<Routes>} in \file{src/app.rs}, not \code{App::route} in
\file{src/main.rs}. And \emph{the client--server call protocol} is a Leptos concern: the
\code{\#[server]} macro generates the endpoint, the serialisation, and the client-side stub for
you. You do not write controllers.

\subsection*{Actix-web 4 --- the HTTP server}

Actix's job in this application is narrow, and keeping it narrow is a design goal. It listens on
a socket, serves static files out of \file{target/site}, hands unmatched requests to Leptos, and
holds shared application state (the connection pool) so that server functions can extract it.
It should not contain business logic and it should not contain page routes. If you find yourself
writing an Actix handler that returns JSON for your own front end to consume, stop --- that is
what a server function is for, and it will give you type safety across the wire that a
hand-written JSON endpoint cannot.

\subsection*{cargo-leptos --- the build orchestrator}

This is the piece most worth understanding concretely, because it is doing more than it appears
to. Version 0.3.6 is installed. When you run \code{cargo leptos watch} it:

\begin{enumerate}
  \item compiles the library target to \code{wasm32-unknown-unknown} with the \code{hydrate}
    feature, then runs \code{wasm-bindgen} over the result;
  \item compiles the binary target natively with the \code{ssr} feature;
  \item compiles \file{style/main.scss} using a \code{sass} binary it downloaded itself;
  \item copies \file{assets/} into \file{target/site/};
  \item starts the server binary, and restarts or reloads on change.
\end{enumerate}

Point three surprises people. \code{sass} is \emph{not} installed on this machine, yet the build
works, because cargo-leptos fetches and caches its own toolchain.

\begin{verified}
\begin{lstlisting}[style=out]
$ command -v sass  ->  (not found)
$ ls ~/.cache/cargo-leptos/
  sass-1.86.0/           tailwindcss-v4.2.1/
  wasm-bindgen-0.2.125/  wasm-bindgen-0.2.126/
\end{lstlisting}
\end{verified}

Note the second entry. A Tailwind~4 CLI has \emph{already} been downloaded into that cache. This
matters for Chapter~\ref{ch:tailwind}: adopting Tailwind here requires no \code{npm}, no
\file{package.json}, and no Node dependency in your build --- cargo-leptos will manage the binary
the same way it manages \code{sass}.

\subsection*{Diesel 2.3 with r2d2 --- typed SQL and connection pooling}

Diesel gives you compile-time-checked queries against the schema in \file{src/schema.rs}. It is
worth being precise about one property that shapes the whole server design: \textbf{Diesel, as
configured here, is synchronous}. Every query blocks the thread that issues it until Postgres
answers. Actix and Leptos server functions are asynchronous. Reconciling those two facts is not
optional and it is not a micro-optimisation --- Chapter~\ref{ch:blocking} explains why a direct
call will damage your server under concurrency, and gives the two-line wrapper that fixes it.

\subsection*{PostgreSQL 17 --- and \code{diesel-cli}}

The database runs in Docker on port 5439. Migrations are plain SQL pairs under
\file{migrations/}, applied by \code{diesel migration run} and reverted by
\code{diesel migration revert}. Chapter~\ref{ch:migrations} is about the discipline that keeps
that second command working, because right now it does not.

\subsection*{Tailwind CSS --- not yet present}

Currently \file{style/main.scss} is three lines of Sass. Chapter~\ref{ch:tailwind} covers the
switch. The short version: Tailwind replaces the Sass file rather than joining it, and the change
is four lines of \file{Cargo.toml} plus one new CSS file.

\section{Where new code goes}

To close the chapter, the practical form of the above. When you are about to write something and
are unsure where it belongs, this table answers it.

\begin{table}[htbp]
\centering
\caption{A decision table for placing new code.}
\label{tab:placement}
\small
\begin{tabular}{@{}L{0.40\textwidth}L{0.54\textwidth}@{}}
\toprule
\textbf{You are writing\ldots} & \textbf{It belongs in\ldots} \\
\midrule
a page or a visual component & \file{src/ui/} (new), reached from \file{src/app.rs} \\
a URL that a human types or bookmarks & \code{<Routes>} in \file{src/app.rs} \\
a call from browser to server & \file{src/<domain>/api.rs} (new), \emph{ungated} \\
a type that travels over the wire & \file{src/<domain>/dto.rs} (new), \emph{ungated} \\
a SQL query & \file{src/<domain>/repository.rs}, gated on \code{ssr} \\
a business rule or invariant & \file{src/<domain>/} beside the model, gated on \code{ssr} \\
an error the UI has to display & \file{src/error.rs} (new), \emph{ungated} \\
process startup and wiring & \file{src/main.rs}, once, before the server starts \\
an Actix route & almost certainly nothing --- reconsider; use a server function \\
\bottomrule
\end{tabular}
\end{table}

The word \emph{ungated} in that table is doing a lot of work, and the reason is the subject of
Chapter~\ref{ch:boundary}. The new files it names are laid out in full in Chapter~\ref{ch:layers}.
